\documentclass[12pt]{article}
\usepackage[utf8]{inputenc}
\usepackage[T1]{fontenc}
\usepackage[margin=1in]{geometry}
\usepackage{setspace}
\usepackage{xcolor}
\usepackage[colorlinks=true, urlcolor=blue!60!black, linkcolor=blue!60!black]{hyperref}
\setlength{\parindent}{0em}
\setlength{\parskip}{0.45em}
\renewcommand{\baselinestretch}{1.0}
\begin{document}
\begin{center}
{\Large \textit{Smith v. Doe}, 538 U.S. 84 (2003)\par}
\vspace{0.3em}
{\normalsize Case Brief --- One-Page Edition\par}
\end{center}

\textbf{Memory Jogger.} Retroactive sex offender registration and notification is civil regulation, not ex post facto punishment (\href{https://www.law.cornell.edu/supct/html/01-729.ZO.html}{Cornell LII}).

\textbf{Facts.} Alaska's 1994 Act required sex offenders, including those convicted earlier, to register and verify information for 15 years or life, and posted most of it online. Two men who pleaded \textit{nolo contendere} to aggravated sexual abuse of a minor before the Act, and completed rehabilitation programs, sued under \S~1983.

\textbf{Procedural History.} The district court upheld the Act; the Ninth Circuit found punitive effects; the Supreme Court reversed.

\textbf{Votes.} 6--3: Kennedy for the Court; Thomas concurred; Souter concurred in the judgment; Stevens and Ginsburg (with Breyer) dissented.

\textbf{Holding.} The legislature intended a civil scheme, and no ``clearest proof'' showed punitive effects.

\textbf{Reasoning.} The Act aimed at ``protecting the public from sex offenders'' through an administrative agency. Under the \textit{Mendoza-Martinez} factors, it is unlike colonial shaming, imposes no physical restraint, rationally serves public safety, and may act categorically. Stevens would treat burdens imposed on all and only convicted offenders as punishment; Ginsburg found the Act excessive for lacking any path off the registry.

\textbf{Examples.} \emph{Same-Side:} civil commitment (\textit{Hendricks}). \emph{Opposite-Side:} harsher later Guidelines (\textit{Peugh}). \emph{Fence-Sitter:} lifetime GPS monitoring.

\textbf{Critique.} The ``clearest proof'' standard invites relabeling punishment as civil, and the ``frightening and high'' recidivism claim rested on a disavowed magazine article.

\textbf{Key Quotation.} ``If the intention of the legislature was to impose punishment, that ends the inquiry. If, however, the intention was to enact a regulatory scheme that is civil and nonpunitive, we must further examine whether the statutory scheme is `so punitive either in purpose or effect as to negate [the State's] intention' to deem it `civil.'\,'' --- Justice Kennedy, 538 U.S. at 92.

\end{document}
